\let\origsectionmark\sectionmark\renewcommand*{\sectionmark}[1]{\origsectionmark{Hohenberg--Kohn theorems and Kohn--Sham construction}}% short running head
\section{Hohenberg--Kohn theorems and the Kohn--Sham construction}\label{app:D.5}
\let\sectionmark\origsectionmark

The ground-state density of the electrons in silicon, a function of three variables, fixes the potential of the nuclei and with it the whole Hamiltonian, every excited state included. This is the first theorem of Hohenberg and Kohn~\cite{hohenberg1964inhomogeneous}. Their second theorem turns the ground-state energy into the minimum of a functional of the density, and Kohn and Sham~\cite{kohn1965self} made that minimum computable with one-electron equations, the Kohn--Sham equations that every band structure in this book comes from (\cref{sec:3.4}). This section proves the two theorems, derives the equations and the total energy, and states what the eigenvalues mean. It uses the constrained-search formulation of Levy~\cite{levy1979universal}, which needs fewer assumptions than the original proof, and follows Ref.~\cite[Chaps.~6 and 7]{martin2004electronic}.

\subsection*{The density determines the potential}

Take $N_e$ electrons in an external potential $U_1(\vect{r})$, the attraction of the clamped nuclei, with
\begin{equation*}
  H = H_{\mathrm{kin}} + H_{ee} + \sum_{i=1}^{N_e}U_1(\vect{r}_i), \qquad H_{\mathrm{kin}} = \sum_i\frac{\vect{p}_i^2}{2m_e} ,
\end{equation*}
and $H_{ee}$ the Coulomb repulsion of \cref{eq:2.3-Hee-density}. Only the last term depends on the material, and it enters every expectation value through the density alone: $\langle\Psi|\sum_iU_1(\vect{r}_i)|\Psi\rangle = \int\dd^3r\,U_1(\vect{r})\,\nd(\vect{r})$, with $\nd(\vect{r})$ the electron density of $\Psi$.

Suppose two potentials $U_1$ and $U_1'$, which differ by more than a constant, had nondegenerate ground states $\Psi$ and $\Psi'$ with the same density $\nd$. The two states are different. If they were equal, subtracting the two Schr\"odinger equations would give $\sum_i[U_1(\vect{r}_i) - U_1'(\vect{r}_i)]\Psi = (E_0 - E_0')\Psi$, so that the sum would be constant wherever $\Psi \neq 0$, which for a ground state is almost everywhere, and the potentials would differ by a constant. Since the ground state of $H$ is nondegenerate, $\Psi'$ has a strictly higher energy in $H$:
\begin{equation*}
  E_0 < \langle\Psi'|H|\Psi'\rangle = \langle\Psi'|H'|\Psi'\rangle + \int\dd^3r\,\bigl[U_1(\vect{r}) - U_1'(\vect{r})\bigr]\nd(\vect{r}) = E_0' + \int\dd^3r\,\bigl[U_1 - U_1'\bigr]\nd .
\end{equation*}
The same argument with the primed and unprimed quantities exchanged gives $E_0' < E_0 + \int\dd^3r\,[U_1' - U_1]\nd$. Adding the two inequalities, the integrals cancel and $E_0 + E_0' < E_0' + E_0$, a contradiction. Hence the ground-state density determines $U_1$ up to a constant~\cite{hohenberg1964inhomogeneous}. It also fixes $N_e = \int\dd^3r\,\nd$, so it fixes $H$, and through $H$ every property of the system. The proof says nothing about how to find any of them from $\nd$.

\subsection*{The variational principle}

For a density $\nd$ define
\begin{equation}
  \mathcal{F}[\nd] = \min_{\Psi\to\nd}\langle\Psi|H_{\mathrm{kin}} + H_{ee}|\Psi\rangle ,
  \label{eq:D.5-levy}
\end{equation}
where the minimum is over all normalized antisymmetric states of $N_e$ electrons that have the density $\nd$~\cite{levy1979universal}. The functional $\mathcal{F}$ is universal: it does not depend on the nuclei, and the same $\mathcal{F}$ serves silicon, germanium and an isolated atom. Every reasonable density, nonnegative, integrating to $N_e$ and with a finite kinetic energy, comes from some determinant, so $\mathcal{F}$ is defined on all of them. For a given potential let
\begin{equation*}
  E_{U}[\nd] = \mathcal{F}[\nd] + \int\dd^3r\,U_1(\vect{r})\,\nd(\vect{r}) .
\end{equation*}
For any $\nd$, the state $\Psi_{\nd}$ that realizes the minimum in \cref{eq:D.5-levy} has $\langle\Psi_{\nd}|H|\Psi_{\nd}\rangle = E_U[\nd]$, and by the ordinary variational principle this is at least $E_0$. The minimum over all states, on the other hand, can be done in two steps, first over the states with a given density and then over the densities:
\begin{equation*}
  E_0 = \min_\Psi\langle\Psi|H|\Psi\rangle = \min_{\nd}\Bigl[\min_{\Psi\to\nd}\langle\Psi|H_{\mathrm{kin}} + H_{ee}|\Psi\rangle + \int\dd^3r\,U_1\nd\Bigr] = \min_{\nd}E_U[\nd] .
\end{equation*}
The ground-state energy is the minimum of $E_U$ over densities, and the minimum is reached at the ground-state density. This is the second theorem of Hohenberg and Kohn, here without the assumption of a nondegenerate ground state, and without the requirement that every trial density be the ground-state density of some potential. Where $\mathcal{F}$ is differentiable, the minimum under the constraint $\int\nd = N_e$ obeys $\delta\mathcal{F}/\delta\nd(\vect{r}) + U_1(\vect{r}) = \mu_{\mathrm{ch}}$, with the chemical potential as the Lagrange multiplier.

\subsection*{The Kohn--Sham equations}

Nobody knows $\mathcal{F}$, and the kinetic energy is its hardest part. Kohn and Sham computed that part exactly for a system of independent electrons with the same density. Define the independent-electron kinetic energy, written $T_{\mathrm{s}}$ in the literature, by the same constrained search over determinants $\Phi$ alone,
\begin{equation*}
  K_{\mathrm{s}}[\nd] = \min_{\Phi\to\nd}\langle\Phi|H_{\mathrm{kin}}|\Phi\rangle ,
\end{equation*}
and the Hartree energy $E_{\mathrm{H}}[\nd] = \tfrac12\int\dd^3r\,\dd^3r'\,e^2\nd(\vect{r})\nd(\vect{r}')/(4\pi|\vect{r} - \vect{r}'|)$, the classical electrostatic energy of the density. The exchange--correlation energy is everything else:
\begin{equation}
  \mathcal{F}[\nd] = K_{\mathrm{s}}[\nd] + E_{\mathrm{H}}[\nd] + E_{\mathrm{xc}}[\nd] .
  \label{eq:D.5-Exc}
\end{equation}
It holds the exchange energy, which removes the self-repulsion contained in $E_{\mathrm{H}}$ as in \cref{app:D.4}, the correlation energy, and the part of the true kinetic energy that independent electrons miss. It is the only part that has to be approximated.

Write the minimizing determinant through its orbitals $\psi_\gamma$, with occupations $n_\gamma$ equal to one for the $N_e$ occupied orbitals and zero otherwise, and the density $\nd(\vect{r}) = \sum_\gamma n_\gamma\sum_s|\psi_\gamma(\vect{r}, s)|^2$. Then
\begin{equation*}
  E_U = \sum_\gamma n_\gamma\int\dd x\,\psi^*_\gamma(x)\Bigl(-\frac{\nabla^2}{2m_e}\Bigr)\psi_\gamma(x) + \int\dd^3r\,U_1\nd + E_{\mathrm{H}}[\nd] + E_{\mathrm{xc}}[\nd] ,
\end{equation*}
which depends on the orbitals explicitly through the kinetic term and implicitly through $\nd$. We vary $\psi^*_\gamma(x)$, with Lagrange multipliers for orthonormality chosen diagonal, as in \cref{app:D.4}. The chain rule gives $\delta E_{\mathrm{H}}/\delta\psi^*_\gamma(x) = n_\gamma v_{\mathrm{H}}(\vect{r})\psi_\gamma(x)$, with the Hartree potential $v_{\mathrm{H}} = \delta E_{\mathrm{H}}/\delta\nd$ of \cref{app:D.4}, and similarly for $E_{\mathrm{xc}}$. Dividing by $n_\gamma$,
\begin{equation}
  \begin{gathered}
  \Bigl[-\frac{\nabla^2}{2m_e} + V_{\mathrm{KS}}(\vect{r})\Bigr]\psi_\gamma(x) = E_\gamma\,\psi_\gamma(x), \\
  V_{\mathrm{KS}}(\vect{r}) = U_1(\vect{r}) + v_{\mathrm{H}}(\vect{r}) + v_{\mathrm{xc}}(\vect{r}), \qquad
  v_{\mathrm{xc}}(\vect{r}) = \frac{\delta E_{\mathrm{xc}}}{\delta\nd(\vect{r})} ,
  \end{gathered}
  \label{eq:D.5-KS}
\end{equation}
with the density built from the $N_e$ orbitals of lowest $E_\gamma$. These are the Kohn--Sham equations~\cite{kohn1965self}. They have the form of the Hartree--Fock equations \cref{eq:D.4-HF}, but the potential $V_{\mathrm{KS}}$ is local, a multiplicative function of $\vect{r}$, where the Fock term was an integral operator. And if $E_{\mathrm{xc}}$ were exact, the density they produce would be the exact ground-state density of the interacting electrons. Filling the lowest orbitals assumes that the minimizing density is the ground-state density of independent electrons in some local potential, the noninteracting $v$-representability of the density; for the crystals of this book it is taken for granted. In a crystal, $\gamma$ becomes the band and crystal momentum $n\vect{k}$, $E_\gamma$ the band energy $\varepsilon_{n\vect{k}}$ and $n_\gamma$ the occupation $f_{n\vect{k}}$, and the density is \cref{eq:E.1-density}. The equations are nonlinear, since $V_{\mathrm{KS}}$ depends on the density, and they are iterated to self-consistency as \cref{app:E.1} describes.

\subsection*{Total energy and the eigenvalues}

Multiplying \cref{eq:D.5-KS} by $n_\gamma\psi^*_\gamma$, integrating and summing gives $\sum_\gamma n_\gamma E_\gamma = K_{\mathrm{s}} + \int\dd^3r\,V_{\mathrm{KS}}\nd$, and $\int v_{\mathrm{H}}\nd = 2E_{\mathrm{H}}$. The ground-state energy is therefore
\begin{equation}
  E_0 = \sum_\gamma n_\gamma E_\gamma - E_{\mathrm{H}}[\nd] - \int\dd^3r\,v_{\mathrm{xc}}(\vect{r})\nd(\vect{r}) + E_{\mathrm{xc}}[\nd] .
  \label{eq:D.5-Etot}
\end{equation}
As in Hartree--Fock, \cref{eq:D.4-Etot}, the sum of the eigenvalues counts the interaction twice and the other terms correct for it.

The eigenvalues are Lagrange multipliers, and \cref{eq:D.5-KS} alone does not make them removal or addition energies. One exact statement connects them to energies. Let the occupations $n_\gamma$ take fractional values and keep the orbitals self-consistent at each set of $n_\gamma$. Then $E_U$ is a function of the occupations, and
\begin{equation}
  \frac{\partial E_U}{\partial n_\gamma} = E_\gamma ,
  \label{eq:D.5-janak}
\end{equation}
which is Janak's theorem~\cite{janak1978proof}. The explicit dependence on $n_\gamma$, through the kinetic term and through $\nd$, gives the expectation value of the Kohn--Sham Hamiltonian in $\psi_\gamma$, which is $E_\gamma$. The implicit dependence through the orbitals gives, by \cref{eq:D.5-KS},
\begin{equation*}
  \sum_{\gamma'}n_{\gamma'}E_{\gamma'}\,\frac{\partial\langle\psi_{\gamma'}|\psi_{\gamma'}\rangle}{\partial n_\gamma} = 0 ,
\end{equation*}
because every orbital stays normalized. For the exact functional the ground-state energy of a fractional number of electrons is a straight line between integers~\cite{perdew1982density}. Integrating \cref{eq:D.5-janak} from $N_e - 1$ to $N_e$ in the highest occupied level, whose eigenvalue is then constant, gives $E_{N_e} - E_{N_e-1} = E_{\mathrm{HOMO}}$, with $E_{\mathrm{HOMO}}$ the eigenvalue of the highest occupied molecular orbital (HOMO): the highest occupied Kohn--Sham eigenvalue is minus the ionization energy $I$. In the same way, the eigenvalue of the level that fills between $N_e$ and $N_e + 1$ is minus the electron affinity $A$. No other eigenvalue has an exact meaning.

The fundamental gap is $E_{\mathrm{g}} = I - A$, the difference of the two removal and addition energies at the edge. It is not the Kohn--Sham gap $E_{\mathrm{g}}^{\mathrm{KS}}$, the difference of the lowest empty and highest occupied eigenvalues of the $N_e$-electron potential. The straight lines on either side of $N_e$ have the slopes $-I$ and $-A$, so the chemical potential jumps at $N_e$ by the fundamental gap. The derivative of the Kohn--Sham kinetic energy with respect to the electron number jumps only by $E_{\mathrm{g}}^{\mathrm{KS}}$, since the added electron enters the lowest unoccupied molecular orbital (LUMO). The Hartree energy is smooth in $\nd$, so the rest of the jump belongs to $E_{\mathrm{xc}}$. As the electron number passes through $N_e$, the derivative $\delta E_{\mathrm{xc}}/\delta\nd$ of the exact functional therefore jumps by a constant $\Delta_{\mathrm{xc}}$, the derivative discontinuity, and every eigenvalue jumps with it; the Hartree and kinetic parts do not jump. The level that fills above $N_e$ therefore has the eigenvalue $E_{\mathrm{LUMO}} + \Delta_{\mathrm{xc}}$, and
\begin{equation}
  E_{\mathrm{g}} = I - A = E_{\mathrm{g}}^{\mathrm{KS}} + \Delta_{\mathrm{xc}}
  \label{eq:D.5-gap}
\end{equation}
for the exact functional~\cite{perdew1983physical,sham1983density}. Even exact Kohn--Sham bands have the wrong gap. The local and semilocal functionals below are smooth functions of the density and have no discontinuity; this is the band-gap problem of \cref{sec:3.4}.

\subsection*{The local-density approximation}

The local-density approximation takes the exchange--correlation energy per electron of the uniform gas, $\varepsilon_{\mathrm{xc}}(\nd)$, at the local density,
\begin{equation}
  E_{\mathrm{xc}}^{\mathrm{LDA}}[\nd] = \int\dd^3r\,\nd(\vect{r})\,\varepsilon_{\mathrm{xc}}\bigl(\nd(\vect{r})\bigr), \qquad
  v_{\mathrm{xc}}^{\mathrm{LDA}}(\vect{r}) = \frac{\dd\bigl[\nd\,\varepsilon_{\mathrm{xc}}(\nd)\bigr]}{\dd\nd}\bigg|_{\nd = \nd(\vect{r})} .
  \label{eq:D.5-LDA}
\end{equation}
It is exact for a uniform density and accurate where the density varies slowly on the scale of the local $1/k_{\mathrm{F}}$~\cite{kohn1965self}. The exchange part is \cref{eq:D.4-Ex}, $\varepsilon_{\mathrm{x}}(\nd) = -3e^2k_{\mathrm{F}}/16\pi^2$ with the local $k_{\mathrm{F}} = (3\pi^2\nd)^{1/3}$, which follows from \cref{eq:2.5-fermi}. Since $\nd\,\varepsilon_{\mathrm{x}} \propto \nd^{4/3}$, the exchange potential is
\begin{equation*}
  v_{\mathrm{x}}^{\mathrm{LDA}} = \frac43\,\varepsilon_{\mathrm{x}} = -\frac{e^2k_{\mathrm{F}}}{4\pi^2} = \Sigma_{\mathrm{x}}(k_{\mathrm{F}}) ,
\end{equation*}
the Hartree--Fock exchange of an electron at the Fermi surface of the gas, \cref{eq:D.4-gas} at $F_{\mathrm{L}}(1) = \tfrac12$. Slater had earlier proposed to replace the Fock term by its average over the occupied states~\cite{slater1951simplification}, which in the gas is $\langle\Sigma_{\mathrm{x}}\rangle = \tfrac34\Sigma_{\mathrm{x}}(0) = \tfrac32v_{\mathrm{x}}^{\mathrm{LDA}}$; Kohn and Sham pointed out that the variational principle leads to two thirds of Slater's potential~\cite{kohn1965self}. At the valence density of silicon, $v_{\mathrm{x}}^{\mathrm{LDA}} = \SI{-8.29}{eV}$ and Slater's average is \SI{-12.4}{eV}. The correlation part of $\varepsilon_{\mathrm{xc}}$ has no closed form. It is taken from quantum Monte Carlo energies of the gas~\cite{ceperley1980ground}, fitted as a function of $r_s$~\cite{perdew1981self}. Unlike the Fock term, the local exchange potential does not cancel the self-repulsion of an electron in $v_{\mathrm{H}}$ exactly, and the residual self-interaction is one source of the error of the local-density approximation in localized states~\cite{perdew1981self}. The generalized-gradient approximations add a dependence on $|\nabla\nd|$, and hybrid functionals replace part of the local exchange by the Fock term built from the Kohn--Sham orbitals; \cref{app:E.1} describes the functionals in use and their accuracy for silicon.
